\section{Discussion and Limitations}
\label{sec:discussion}
In this section, we discuss the broader implications of our diagnostic protocol for educational technology research and practice. We then define the scope boundaries of our work and address the primary limitations of the current study.

\subsection{Implications}
The proposed protocol encourages KT evaluation to move beyond aggregate AUC. 
For researchers, per-stratum reporting combined with reliability flags makes claims about sparse or cold-start improvements verifiable and auditable. 
For developers of educational systems, recognizing that calibration can degrade on sparse KCs implies that any threshold-based interventions (e.g., mastery thresholds) require per-stratum calibration validation rather than relying solely on overall accuracy.
For reviewers and the broader community, the fact that our L8 sanity check successfully caught a subtle prediction--label misalignment highlights the critical need for routine, standardized leakage and predictive-sanity audits in all KT sequence modeling workflows.

\subsection{Scope Boundaries}
This paper is intentionally limited to evaluation diagnostics. It does not introduce a new KT model, a self-supervised learning (SSL) objective, a graph neural network (GNN), a graph augmentation strategy, a learning path recommender (LPR), or a distillation method. Those directions require separate methodological papers and should be evaluated using the protocol introduced here. A complementary line of work in our group addresses leakage control in concept-graph construction for graph-augmented KT. These two protocols are designed to be composable in future structure-aware KT studies.

\subsection{Threats to Validity and Limitations}
\label{sec:threats}
We analyze the threats to the validity of our diagnostic protocol across standard dimensions:
\begin{itemize}
    \item \textbf{Internal Validity:} Risks to internal validity include potential preprocessing bugs, split leakage, sequence--label alignment errors, and implementation-specific training differences across the baselines. We mitigate these risks through a standardized interaction data schema, rigid validation-fold early stopping, and our automated leakage audit checklist. Additionally, we note that our classical baseline (IRT) under learner-based splits acts as a base-rate reference due to the lack of ability estimates for unseen learners; developing a classical baseline that effectively exploits item difficulty for unseen learners would strengthen future comparative work.
    \item \textbf{Temporal Alignment Validity:} Sequence modeling in KT is notoriously prone to ``silent bugs'' such as off-by-one prediction--label misalignments. Our protocol actively defends against this via the warm-cohort predictive sanity audit (L8). During validation, this automated check successfully flagged an underlying library alignment discrepancy that otherwise would have systematically corrupted the evaluation.
    \item \textbf{Construct Validity:} Construct validity refers to whether our metrics accurately capture latent learner mastery and concept calibration. In KT, KC definitions depend heavily on dataset annotation quality, Q-matrix provenance, and how multi-skill items are processed. For example, expert-tagged KCs in ASSISTments 2012 and Junyi Academy may introduce varying granularities, which directly impact strata boundaries.
    \item \textbf{External Validity:} The pronounced heterogeneity observed across datasets---such as the clear calibration degradation gradient in ASSISTments 2012, the inverted calibration patterns in XES3G5M, and the entirely empty very-sparse stratum in Junyi Academy---supports the necessity of the protocol. Models cannot be assumed to behave consistently across different sparsity distributions, making dataset-local validation essential.
    \item \textbf{Statistical Conclusion Validity:} Statistical conclusion validity concerns the reliability of our metric estimates. Very sparse buckets contain small samples, meaning strata differences must be interpreted in conjunction with explicit sample size counts and reliability flags. Additionally, our temporal diagnostics rely on a single corrected run with a fixed chronological cutoff, and our Brier decomposition utilizes a binned empirical approximation.
\end{itemize}
